\documentclass[10pt,a4paper]{amsart}
\usepackage[margin=19mm]{geometry}
\usepackage{amsmath,amssymb,mathtools}
\usepackage[scr=rsfs,cal=euler]{mathalfa}
\usepackage{xfrac}
\usepackage[unicode,hidelinks,bookmarks=false]{hyperref}
\newcommand{\ie}{i.e.\ }
\newcommand{\cf}{cf.\ }
\newcommand{\resp}{respectively\ }
\newcommand{\Subsection}{\subsection}
\providecommand{\nobreakdash}{\nobreak\hbox{-}}
\setlength{\emergencystretch}{2em}
\multlinegap0pt
\usepackage{amssymb}
\usepackage{float}
\usepackage{mathtools}
\usepackage{siunitx}
\usepackage{algorithm}\usepackage{algpseudocode}
\newtheorem{theorem}{Theorem}
\newcommand{\E}{\mathcal{F}}
\newcommand{\X}{\mathcal{X}}
\newcommand{\argmin}{\mathop{\mathrm{argmin}}}
\newcommand{\grad}{\nabla}%
\newcommand{\norm}[1]{\left\| #1 \right\|}
\newcommand{\ph}{\phi}%
\title{High-Order Variable-Scaled Energetic Variational Neural Networks for Phase-Field Gradient Flows}
\author{Xiaobo Jing, Leiyi Dong, Jia Zhao}
\date{Source: arXiv:2609.25489v1 (CC BY 4.0)}
\begin{document}
\maketitle

\noindent\emph{Source and licence.} High-Order Variable-Scaled Energetic Variational Neural Networks for Phase-Field Gradient Flows. Version arXiv:2609.25489v1 (CC BY 4.0), distributed by arXiv under the Creative Commons Attribution 4.0 International licence, http://creativecommons.org/licenses/by/4.0/, as recorded in the arXiv licence metadata for this exact version and shown on the abstract page https://arxiv.org/abs/2609.25489v1. Full licence notice: \texttt{licenses/arxiv-2609.25489-CC-BY-4.0.txt}.\par\smallskip

\noindent\textit{Editorial scope.} Counted unit: the article's theorem thm:first\_order (source0.tex lines 272--301). The article's complete original proof of it is delivered with it (source0.tex lines 303--362, one \texttt{proof} environment of 60 lines). No mathematics is written, repaired or re-derived: the statement and the proof are the article's own lines, in the article's own notation, and no retained line is substituted or re-flowed.  No further source-local context is needed: every cross-reference of every retained line resolves inside the two delivered spans.  Nothing was omitted from any retained span, so the package carries no omission record.  No retained line contains a citation, so the wrapper carries no bibliography and no bibliography entry.  No retained line contains a figure environment, an image-inclusion command or drawing code, so no image is delivered and the package declares no asset dependency.  Editorial formatting only, in addition: an AMS \texttt{amsart} preamble that restates the article's own declarations and macros used by the retained lines, namely the environments \texttt{theorem} on separate counters, together with the standard packages those lines need.  The retained environments print in this document as Theorem 1 in order of appearance, whereas the article prints the same items with the numbering of its own full document.  The complete native source of the article as distributed in the arXiv e-print for this exact version is bundled unchanged as \texttt{sources/211307/source0.tex}.
\begin{theorem}[First-order temporal convergence for the discrete gradient system]
\label{thm:first_order}
Let $\X_h\simeq\mathbb{R}^N$ be equipped with a discrete inner product, and let $\E_h\in C^1(\X_h)$ be the corresponding discrete energy. Let $\ph_h\in C^2([0,T];\X_h)$ solve
\[
\partial_t\phi_h(t)=-\grad\E_h(\phi_h(t)).
\]
Assume that $\grad\E_h$ is Lipschitz continuous on a set $B_h\subset\X_h$ that contains the exact trajectory $\{\phi_h(t):0\le t\le T\}$ and the discrete iterates $\{\phi_h^n\}$,
\[
\norm{\grad\E_h(u)-\grad\E_h(w)} \le L_h\,\norm{u-w}
\qquad \text{for all } u,w\in B_h.
\]
Let $\{\phi_h^n\}$ be generated from $\phi_h^0=\phi_h(0)$ by the discrete minimizing-movements step
\[
\phi_h^{n+1}\in
\argmin_{\phi\in\X_h}
\left\{
\frac{1}{2\tau}\norm{\phi-\phi_h^n}^2+\E_h(\phi)
\right\},
\]
assuming that a minimizer exists at every step. Then, for every $\tau>0$ satisfying $\tau L_h\le 1/2$,
\begin{equation}
\max_{0\le n\le \lfloor T/\tau\rfloor}\,
\norm{\phi_h(t_n)-\phi_h^n}
\le
\frac{M_{2,h}\,T\,e^{2L_hT}}{2}\,\tau,
\qquad
M_{2,h}:=\sup_{t\in[0,T]}\norm{\partial_{tt}\phi_h(t)} .
\label{eq:first-order-bound}
\end{equation}
\end{theorem}

\begin{proof}

The Euler--Lagrange equation of the discrete step,
\begin{equation}
\frac{\phi_h^{n+1}-\phi_h^n}{\tau}
+\grad\E_h(\phi_h^{n+1})=0,
\label{eq:euler-lagrange}
\end{equation}
identifies the scheme as implicit Euler.

Since $\phi_h^{n+1}$ minimizes the differentiable discrete functional over $\X_h$, it is a critical point, and the Euler--Lagrange equation \eqref{eq:euler-lagrange} holds exactly.

A Taylor expansion with integral remainder about $t_{n+1}$ gives
\[
\phi_h(t_n)=\phi_h(t_{n+1})-\tau\,\partial_t\phi_h(t_{n+1})+R_h^n,
\qquad
\norm{R_h^n}\le\frac{\tau^2}{2}\,M_{2,h},
\]
and substituting $\partial_t\phi_h(t_{n+1})=-\grad\E_h(\phi_h(t_{n+1}))$ yields
\[
\phi_h(t_{n+1})
=
\phi_h(t_n)-\tau\,\grad\E_h(\phi_h(t_{n+1}))-R_h^n .
\]

Set $\xi^n:=\phi_h(t_n)-\phi_h^n$. Subtracting \eqref{eq:euler-lagrange} from the identity above gives
\[
\xi^{n+1}
=
\xi^n
-\tau\Bigl(\grad\E_h(\phi_h(t_{n+1}))
-\grad\E_h(\phi_h^{n+1})\Bigr)
-R_h^n .
\]
The Lipschitz bound on $B_h$ then gives
\[
\norm{\xi^{n+1}}
\le
\norm{\xi^n}+\tau L_h\norm{\xi^{n+1}}
+\frac{\tau^2}{2}M_{2,h}.
\]
For $\tau L_h\le 1/2$, we have $(1-\tau L_h)^{-1}\le 1+2\tau L_h\le e^{2\tau L_h}$, hence
\[
\norm{\xi^{n+1}}
\le
e^{2\tau L_h}
\Bigl(\norm{\xi^n}+\frac{\tau^2}{2}M_{2,h}\Bigr).
\]
Since $\xi^0=0$, iterating gives, for $n\tau\le T$,
\[
\norm{\xi^n}
\le
\frac{\tau^2}{2}M_{2,h}\sum_{k=1}^{n} e^{2k\tau L_h}
\le
\frac{\tau^2}{2}M_{2,h}\,n\,e^{2L_hT}
\le
\frac{M_{2,h}\,T\,e^{2L_hT}}{2}\,\tau,
\]
which is \eqref{eq:first-order-bound}.
\end{proof}

\end{document}
